% !TeX program = xelatex
% 单文件入口，使用 XeLaTeX 连续编译两次。
\documentclass[UTF8,11pt,a4paper,fontset=none]{ctexart}
\usepackage[margin=2.2cm]{geometry}
\usepackage{fontspec}
\IfFontExistsTF{FandolSong-Regular}{
  \setCJKmainfont{FandolSong-Regular}[BoldFont=FandolHei-Regular]
  \setCJKsansfont{FandolHei-Regular}
  \setCJKmonofont{FandolSong-Regular}
}{
  \setCJKmainfont{SimSun}[BoldFont=SimHei]
  \setCJKsansfont{SimHei}
  \setCJKmonofont{SimSun}
}
\usepackage{amsmath,amssymb,booktabs,longtable,array,enumitem,xcolor,listings}
\usepackage{hyperref}
\hypersetup{colorlinks=true,linkcolor=black,urlcolor=blue!55!black,pdftitle={CCF CSP 前三题总讲义编写方案}}
\setlength{\parskip}{0.35em}
\setlength{\emergencystretch}{2em}
\pagestyle{plain}
\setlength{\tabcolsep}{5pt}
\renewcommand{\arraystretch}{1.22}
\setlist{nosep,topsep=0.3em,leftmargin=2em}
\ctexset{section={format=\Large\bfseries},subsection={format=\large\bfseries},subsubsection={format=\normalsize\bfseries}}
\newcolumntype{P}[1]{>{\raggedright\arraybackslash}p{#1}}
\newcommand{\src}[2]{\textbf{来源：}#1，第 #2 页。}
\newcommand{\prob}[1]{\href{https://www.luogu.com.cn/problem/#1}{\texttt{#1}}}
\newcommand{\topic}[4]{\subsection{#1}\textbf{学习目标：}#2\par\textbf{二级目录：}#3\par\textbf{编写与练习安排：}#4\par}
\lstset{language=C++,basicstyle=\ttfamily\small,breaklines=true,columns=fullflexible,keepspaces=true,showstringspaces=false,frame=single,rulecolor=\color{black!25},numbers=left,numberstyle=\tiny,numbersep=7pt,xleftmargin=1.5em}
\title{\textbf{CCF CSP 前三题总讲义编写方案}}
\author{资料汇总\quad 章节规划\quad 训练安排}
\date{2026年10月9日}
\begin{document}
\maketitle

\section*{整理结论}
这套资料适合整理成面向已有 C/C++ 基础学生的 CCF-CSP 认证前三题讲义。教学主线建议确定为：第一题的准确实现、第二题的常用算法与建模、第三题的字符串处理与综合模拟。讲义正文按知识依赖组织，每章再标明与前三题的联系。

现有 7 份资料包含 160 页课件和 7 页 PDF。两份动态规划课件共 71 页，占课件页数约 44.4\%。基础算法素材较充足，第三题的完整案例和程序组织方法需要重点补写。背包高级变式、三分、模拟退火等内容适合放到选读附录。

\textbf{建议成品：}12 章正文加附录，先完成约 160--200 页的主线讲义，另设可按需扩充的选读部分。此页数是编写预算，不是现有材料转换后的实测页数。正文保留可编译代码、推导、样例演算、分层练习和题目来源。

\textbf{本文件的用途：}作为总讲义的编写蓝图，给出资料清单、细化目录、来源对应、勘误、题目入口和授课安排。完整题解、全部代码及第三题案例仍需按本方案逐章编写。本方案中的课时、优先级和篇幅均为教学建议，不代表考试命题或得分保证。

\tableofcontents
\newpage

\section{原始资料清单与适用范围}
\begin{longtable}{P{0.7cm}P{4.5cm}P{1cm}P{8.6cm}}
\toprule 编号 & 文件 & 页数 & 内容及整理用途 \\\midrule\endhead
S1 & 历年前三题考点梳理.pdf & 5 & 第一题 38 条、第二题 38 条、第三题 19 条考点记录；适合做题型导读和选题依据，需补齐题名与认证编号。\\
S2 & 练习题与真题.pdf & 2 & 第一题 8 个、第二题 15 个题目链接，以及 5 个专题题单入口；用作训练索引。\\
S3 & stl.pptx & 42 & 排序、去重、二分接口、离散化、容器、字符串、优先队列、可删堆；拆入工具章和字符串章。\\
S4 & 前缀和、差分、二分、双指针.pptx & 34 & 二分及二分答案、前缀和与差分、双指针；主体适合正文，部分优化算法移至选读。\\
S5 & 搜索、剪枝、暴力.pptx & 13 & DFS、BFS、记忆化搜索、剪枝的题目入口；概念与完整推导需补写。\\
S6 & dp入门.pptx & 18 & 递推、状态设计、最长上升子序列、最大子段和、记忆化、数字三角形、滚动数组。\\
S7 & 动态规划之背包问题.pptx & 53 & 基础背包、分组与混合、多维背包、计数及较多竞赛变式；按必修与选读重新分配。\\
\bottomrule
\end{longtable}

\subsection{考点记录的使用边界}
S1 的条目是资料中的顺序编号，不能直接当成认证编号。它没有为每个条目提供年份、题名和官方题号；第三题仅列出 19 条，与前两题的 38 条并不对齐。因此，不宜据此计算历年命题频率，或推断完整的难度变化趋势。

S2 保存了题目入口，尚未形成“题名、认证编号、主考点、练习层级、解法与讲义页码”的完整索引。不同平台链接也可能对应同一道真题，不能把 23 个入口直接认定为 23 道不重复真题。成稿前应按认证编号去重。

\subsection{适合直接吸收与必须加工的内容}
\begin{itemize}
\item 可改写吸收：前缀和、差分、二分答案、基础 DP 与背包的概念和示例。
\item 需要重新组织：STL 接口列表、散落的题号和作业评讲；改成按问题需求选择工具。
\item 需要转写：大量题面和代码以截图保存，应改成可搜索文字和可编译代码。
\item 需要补写：复杂度估计、输入细节、简单数学、贪心证明、第三题解析器与系统化调试。
\item 需要补齐来源：S1 中未标题名的考点、S5 中“第37次模拟模板展开”等简称。
\end{itemize}

\section{总讲义结构与篇幅预算}
\begin{longtable}{P{0.8cm}P{4.6cm}P{1.6cm}P{7.8cm}}
\toprule 章 & 章名 & 页数建议 & 章节作用 \\\midrule\endhead
01 & 备考方法与程序基本功 & 6--8 & 复杂度、数据范围、输入输出、边界和调试。\\
02 & 基础枚举与小模拟 & 10--12 & 第一题主线，保证规则转写准确。\\
03 & 排序与常用 STL 工具 & 12--14 & 排序、容器选择、统计、离散化。\\
04 & 字符串处理与解析基础 & 12--14 & 为第三题提供输入、扫描和分词能力。\\
05 & 前缀和与差分 & 12--14 & 一维、二维模型与边界处理。\\
06 & 二分与双指针 & 14--16 & 单调性、答案判定、窗口和重复计数。\\
07 & 贪心与基础数学 & 8--10 & 排序决策、交换思想、质因数与简单矩阵运算。\\
08 & 搜索与剪枝 & 10--12 & DFS、BFS、状态空间、正确剪枝。\\
09 & 动态规划入门 & 12--14 & 状态设计、初始化、转移顺序和空间优化。\\
10 & 基础背包与计数 & 10--12 & 01、完全、基础分组背包及语义比较。\\
11 & 第三题综合模拟专题 & 24--30 & 长题意拆解、解析器、树形数据和综合案例。\\
12 & 真题训练与限时复盘 & 16--20 & 专题迁移、分阶段测试、错因归类。\\
\bottomrule
\end{longtable}
上述正文预算合计 146--176 页；补充约 14--24 页的工具速查、题目索引和习题提示，可形成约 160--200 页的主线讲义。页数随完整题面和题解长度调整。

\subsection{阅读层级}
\textbf{必修层：}第 1--6 章、简单贪心、基础搜索、线性 DP、01 与完全背包、第三题模拟方法。\textbf{巩固层：}二维差分、分组背包、递归解析、较复杂状态模拟。\textbf{选读层：}多重背包二进制拆分、混合和多维背包、背包撤销、可删堆、进阶 DP 优化、三分与启发式搜索。

层级按本讲义目标划分，不按原资料所在文件划分。例如，背包课件中的排序贪心题应移至第 7 章；STL 中的字符串应放在第 4 章；搜索课件里的记忆化应与第 9 章统一解释。

\section{十二章详细编写提纲}
\topic{第1章 备考方法与程序基本功}
{能根据数据范围判断解法是否可行，并把题意转换为可检查的输入、输出和处理过程。}
{1.1 前三题训练目标；1.2 时间与空间复杂度；1.3 整数范围与浮点误差；1.4 数组、下标与区间约定；1.5 标准输入输出；1.6 样例外测试与错因记录。}
{以 S1 为导读，补写主体。用同一统计问题比较枚举、排序和映射方案。讲清乘法前提升到长整型、整行输入中的空格、输出格式，以及空集合、单元素和最大规模测试。章末安排一次基础诊断，不将平台得分直接等同于知识掌握。}

\topic{第2章 基础枚举与小模拟}
{能够把文字规则逐条落实，并能说明每个变量表示哪个时刻的状态。}
{2.1 计数与统计；2.2 数位与字符串枚举；2.3 数组和网格模拟；2.4 矩形交与简单几何；2.5 日期及边界；2.6 多规则处理与同步更新；2.7 第一题复盘。}
{用 S1 第一题考点挑选代表题，从 S2 的 8 个第一题入口形成训练池。先给状态表，再推演样例，最后实现。补写“同一轮同时变化时不能直接覆盖旧状态”的反例。正式编写时每个代表题补上题名和认证编号。}

\topic{第3章 排序与常用STL工具}
{能够根据访问、插入、删除、查找和有序性需求选择容器。}
{3.1 多关键字排序与比较器；3.2 vector 与 pair；3.3 去重与离散化；3.4 map 与 set；3.5 queue 与 stack；3.6 priority\_queue；3.7 容器复杂度和常见误用。}
{重组 S3 第 3--19、27--40 页。以计数、排名和动态取最值串起接口。示例应解释迭代器边界、erase 返回值和 map 的默认插入。\prob{P1177}、\prob{P1059}、\prob{P1097} 可作基础训练；\prob{P1090} 留作与贪心的衔接。可删堆移到附录。}

\topic{第4章 字符串处理与解析基础}
{能处理包含空格、分隔符、引号和嵌套结构的输入，并把解析和业务处理分开。}
{4.1 string 与整行输入；4.2 substr、find 与 npos；4.3 扫描指针与分词；4.4 数字解析；4.5 括号、引号与转义；4.6 简单表达式；4.7 字符串键与映射查询。}
{吸收 S3 第 20--26 页及 S1 的字符串补充材料。补写一份逐字符扫描例题和一份按分隔符分词例题，说明为什么遇到引号或嵌套时简单切分不够。C 字符数组作为对照，正文优先使用一致的 C++ 字符串接口。字符串匹配算法放选读，先判断直接匹配能否满足规模。}

\topic{第5章 前缀和与差分}
{能区分静态区间查询、批量区间更新和需要在线处理的操作。}
{5.1 从重复求和到预处理；5.2 一维前缀和；5.3 二维前缀和与容斥；5.4 一维差分；5.5 二维差分；5.6 离线处理的适用条件。}
{主体来自 S4 第 19--25 页。用 \prob{P8218}、\prob{P2367} 分别引入查询与修改，\prob{P13787} 作为二维更新题入口。补写矩形边界图、单点和整段演算。固定一套下标约定，说明前缀数组第 0 行/列和差分的 r+1 位置。高维前缀和转入附录。}

\topic{第6章 二分与双指针}
{能给出单调性理由，明确搜索边界，并说明双指针为什么总移动次数为线性级别。}
{6.1 有序数组查找；6.2 第一个满足与最后一个满足；6.3 二分答案与判定函数；6.4 实数二分选读；6.5 同向与相向双指针；6.6 滑动窗口；6.7 重复元素计数。}
{采用 S4 第 2--14、27--34 页，衔接 S3 的 lower\_bound 和 upper\_bound。代表题为 \prob{P2249}、\prob{P2440}、\prob{P2678}、\prob{P1102}，另保留 LeetCode 713 作为窗口模型。强调窗口乘积题中的正整数前提和 k 小于等于 1 的处理。分数规划、三分和模拟退火移出必修。}

\topic{第7章 贪心与基础数学}
{能说明局部选择为何不会破坏最优解，并掌握支撑第二题的常见数学操作。}
{7.1 排序后的贪心决策；7.2 交换论证与反例；7.3 堆与合并代价；7.4 中位数与距离；7.5 整除与质因数分解；7.6 简单矩阵运算与性质化简。}
{S1 有相关考点，缺少独立成体系讲解，应补写。以 \prob{P1090} 和 S7 第 22--23 页的题目作候选案例，后者作为提高训练。强调“看起来合理”的选择不等于已证明。基础数学只补与现有题型直接相关的内容，暂不扩张为完整数论课程。}

\topic{第8章 搜索与剪枝}
{能从问题中提取状态、转移和终止条件，并区分路径内标记、全局访问和最优距离。}
{8.1 枚举与搜索树；8.2 DFS 和回溯；8.3 BFS 与等权最短步数；8.4 网格连通性；8.5 可行性及最优性剪枝；8.6 搜索与DP的联系。}
{基于 S5 全部内容扩写。\prob{P1706} 引入枚举，\prob{P1219} 展示约束剪枝，\prob{P1443}、\prob{P1135} 讲 BFS，\prob{P1162} 讲连通块。\prob{P3956} 保留为需要仔细设计状态与代价的提高题，不套用普通等权 BFS。\prob{P1120}、\prob{P1514} 放选读；记忆化的完整原理统一在第 9 章讲解。}

\topic{第9章 动态规划入门}
{能够独立写出状态含义、转移、初值、计算顺序和答案位置。}
{9.1 递推与子问题；9.2 状态必须记录什么；9.3 最长上升子序列；9.4 最大子段和；9.5 数字三角形；9.6 二维DP与滚动数组；9.7 记忆化搜索。}
{采用 S6 第 3--16 页，删去重复引入。基础题包括截图中的 \prob{B3637} 与 \prob{P1115}。数字三角形要补边界与负数情形。S6 第 15 页的 \prob{P1439} 在原课件中只要求基础算法获取部分分，应明确标注，不能作为 O(n²) 解法满分练习。\prob{P1541} 等较多维状态放提高层。}

\topic{第10章 基础背包与计数}
{能根据“最多/恰好”“一次/无限次”“最大值/方案数”选择正确的状态与循环。}
{10.1 01背包二维状态；10.2 一维压缩与逆序；10.3 完全背包与正序；10.4 恰好装满与不可达状态；10.5 计数背包；10.6 基础分组背包；10.7 模型对照。}
{优先使用 S7 第 3--12、24--26 页，以 \prob{P1048}、\prob{P1616}、\prob{P1757} 串联。\prob{P1077}、\prob{P1474} 作计数练习候选。多重、混合、多维、撤销和第 k 优等变式进入选读。补一张状态语义对照表，解释循环交换在计数中可能改变组合与排列的含义。}

\topic{第11章 第三题综合模拟专题}
{能把长题意拆成多个可独立检查的部分，完成较长程序并定位错误。}
{11.1 规则清单与状态表；11.2 输入层和执行层；11.3 路径与文本处理；11.4 层级结构与查询；11.5 表达式和递归解析；11.6 对象管理与事件模拟；11.7 完整案例复盘。}
{S1 第 4--5 页仅提供题型线索，S3 提供工具，正文需要重点新写。拟选择 4--6 个完整案例，覆盖字符串扫描、map 查询、树形结构和多规则事件处理。每例展示“题意表格化、对象设计、函数划分、局部测试、整题联调”。涉及“模板展开”的案例须先确认认证编号及题面。}

\topic{第12章 真题训练与限时复盘}
{能够在限定时间内选择解题顺序，保留可验证的基础解法，并从错误中提炼可复用检查项。}
{12.1 第一题稳定性训练；12.2 第二题专题迁移；12.3 第三题渐进实现；12.4 综合限时练习；12.5 错因分类；12.6 二次独立重做。}
{将 S2 的入口重新标注后编入训练池。训练记录包括题号、用时、独立完成程度、错误类型和复做结果。按语法/边界、模型、复杂度、输入解析、状态维护分类复盘。具体取分策略依据已核对题目的子任务设置，不虚构统一的评分规则。}

\section{资料迁移与去重规则}
\begin{longtable}{P{3.8cm}P{2.8cm}P{8.2cm}}
\toprule 原资料位置 & 目标位置 & 处理方式 \\\midrule\endhead
S1 第1--2页 & 第1--2章 & 保留题型总结，补题名与编号；C字符串示例需勘误。\\
S1 第3--4页 & 第5--10章 & 按算法归类，给每条真题建立来源标识。\\
S1 第4--5页 & 第11章 & 用于选案例；缺题面说明的记录先核对再纳入。\\
S2 全部 & 第12章与索引 & 提取文字链接和PDF链接注释，按认证编号去重。\\
S3 第3--8页 & 第3与6章 & 排序、去重、离散化放第3章；二分接口在第6章引用。\\
S3 第9--19、27--40页 & 第3章 & 以问题需求组织容器；合并果子交叉引用第7章。\\
S3 第20--26页 & 第4章 & 完整改写string部分，新增解析例题。\\
S3 第41--42页 & 选读附录 & 可删堆不占必修正文。\\
S4 第2--14、27--34页 & 第6章 & 二分、窗口分别展开；同一道题的多解只保留一份题面。\\
S4 第19--25页 & 第5章 & 重画边界示意，重写可复制代码。\\
S4 第15--18、26页 & 选读附录 & 分数规划、三分、启发式搜索、高维前缀和。\\
S5 全部 & 第8--9章 & BFS、DFS与记忆化重新归类，增加状态设计过程。\\
S6 第3--16页 & 第9章 & 以完整DP五要素组织，滚动数组只讲一次。\\
S6 第17--18页 & 选读附录 & DP优化动机保留，避免在入门阶段展开高级优化。\\
S7 第3--12、24--26页 & 第10章 & 01、完全、分组背包为主体。\\
S7 第22--23页 & 第7章 & 将贪心类作业从背包体系中移出。\\
S7 其余例题与变式 & 选读或训练池 & 按模型细分，避免高级变式挤占综合模拟篇幅。\\
\bottomrule
\end{longtable}

同一知识点只保留一个正式定义和一个基础模板，其他章节使用交叉引用。同一题只保留一份标准题面，多解集中在同一题解中比较。题目包含多个方法时，记录一个主考点和若干辅助考点，避免重复计入训练量。

\section{正式编写前的重点勘误}
\subsection{概念与代码表述}
\begin{enumerate}
\item \textbf{sort 的实现保证。}S3 第4页写成“使用快速排序算法”。应改为标准要求的排序效果及 O(n log n) 比较次数保证，不把具体实现固定成普通快速排序。核对来源：\href{https://eel.is/c++draft/sort}{C++工作草案 sort}。
\item \textbf{unique 的语义。}S3 第5页“把重复数挪到数组后面”会误导读者。应解释为消除相邻等价元素，返回新逻辑末尾；容器尺寸不会自动改变，尾段也不应作为有效结果使用。全局去重通常先排序，再配合 erase。核对来源：\href{https://eel.is/c++draft/alg.unique}{C++工作草案 unique}。
\item \textbf{C字符串赋值与比较。}S1 第2--3页示例应区分字符数组与指针。字符数组在定义后不能用赋值表达式整体赋值，例如 \texttt{S1="AABBCC"}；\texttt{strcmp} 的结果应按正、零、负判断，不能保证为 1、0、-1。对于 \texttt{"aabb"} 和 \texttt{"AABB"}，原示例给出的符号也应复查，示例应明确所用字符编码。
\item \textbf{无后效性的含义。}S6 第7页“当前决策不会对之后的决策产生影响”不准确。应写成：状态已包含预测后续所需的信息，给定状态后不再需要区分到达该状态的不同历史。当前选择当然可以改变后续状态。
\item \textbf{DP初始化。}S6 的楼梯例题需明确 f[0]=1，并约定负下标状态为0或从合法前驱转移。数字三角形的不可达状态和边界不能与合法值混淆，尤其要检查全负数输入。
\item \textbf{背包状态语义。}S7 第4页用“体积和为v”，应与第8页区分“恰好”和“不超过”的初始化对应。讲义每个模板都同时写明状态、初值及最终答案位置，不在两种语义间切换。
\end{enumerate}

\subsection{题意和方法归类}
\begin{enumerate}[resume]
\item \textbf{跳石头判定函数。}S4 第13页写“最多能移走的石头数量”。为了判定距离下界是否可行，应统计满足该距离要求\textbf{至少需要移走}的石头数，并与允许移走数量比较；补上终点处理和单调性说明。
\item \textbf{消失之物。}S7 第32页将 \prob{P4141} 概述为求最大价值，但原题要求去掉每件物品后，恰好装满各容量的\textbf{方案数末位数字}。第34页也转向计数，与前文矛盾。应按计数问题重写整段，不能保留最大价值版本。题面已核对：\href{https://www.luogu.com.cn/problem/P4141}{洛谷P4141}。
\item \textbf{马的遍历。}S5 第5--6页将 \prob{P1443} 放在“记忆化搜索”附近。原题求等代价移动的最少步数，正文应放入 BFS，避免以普通记忆化递归处理有环状态空间。题面已核对：\href{https://www.luogu.com.cn/problem/P1443}{洛谷P1443}。
\item \textbf{最长公共子序列练习。}S6 第15页注明“得到50分即可”。其截图是 \prob{P1439}，正式讲义必须保留“部分分练习”的标记，另选规模匹配的二维DP基础训练，或在选读中补高效解法。
\item \textbf{课程题型范围。}S4 的“可能会考”“不会考”属于原授课备注，整理时换成“本讲义必修/选读”的教学安排。S1、S2 的趋势判断也应限定在已核对的样本范围内。
\end{enumerate}

以上为对可提取文字和关键截图的重点复核，不代替后续全部代码的逐题测试。原始资料保留原样，修订落实到讲义正文，并在编写记录中标明来源页码。

\section{练习体系与现有题目入口}
\subsection{三层练习的设计}
每章配置约 2 道课堂例题、2--3 道基础练习、1--2 道迁移练习，额外提高题进入选读。数量是编写目标，不是现有完整题解数量。基础层要求能够解释代码；迁移层要求独立识别模型；提高层服务于学有余力的学生。

\begin{longtable}{P{2.5cm}P{7.2cm}P{5.1cm}}
\toprule 专题 & 原资料中的候选题 & 编入方式 \\\midrule\endhead
排序和统计 & \prob{P1177}、\prob{P1059}、\prob{P1097}、\prob{P4305} & 第3章基础与容器比较。\\
容器提高 & \prob{P3370}、\prob{P2814}、\prob{P5250} & 先核对完整题面与所需模型，避免只为练接口强行限定解法。\\
二分和窗口 & \prob{P2249}、\prob{P2440}、\prob{P2678}、\prob{P1102}、LeetCode 713 & 第6章从模板到判定建模。\\
前缀与差分 & \prob{P8218}、\prob{P2367}、\prob{P13787} & 第5章一维与二维对照。\\
搜索 & \prob{P1706}、\prob{P1219}、\prob{P1443}、\prob{P1135}、\prob{P1162} & 第8章基础；\prob{P3956} 作提高。\\
线性DP & \prob{B3637}、\prob{P1115}、数字三角形 & 第9章核心例题，完整说明状态和初值。\\
DP迁移 & \prob{P1077}、\prob{P1233}、\prob{P1095}、\prob{P1541}、\prob{P1439} & 根据规模与先修要求分层，P1439标注部分分。\\
基础背包 & \prob{P1048}、\prob{P1616}、\prob{P1757}、\prob{P1474} & 第10章01、完全、分组及计数。\\
背包选读 & \prob{P1776}、\prob{P1833}、\prob{P1507}、\prob{P1064}、\prob{P2732}、\prob{P4141} & 与主线分开，P4141先完成题意勘误。\\
贪心 & \prob{P1090}、\prob{B4071}、\prob{P11376} & 合并果子作基础，后两题作提高候选。\\
\bottomrule
\end{longtable}

其他已出现的选读题号可继续保留在题库：\prob{P8658}、\prob{P1514}、\prob{P1120}、\prob{P1024}、\prob{P2340}、\prob{P5365}、\prob{P1853}、\prob{B4070}、\prob{P5322}、\prob{P1385}、\prob{P2066}、\prob{P2851}、\prob{P5020}、\prob{P1858}、\prob{P3985}、CF19B。S1 的 T219653 是多关键字排序练习入口；是否仍可访问需在选入课堂前确认。

\subsection{精选真题入口}
下表保留 S2 的原始链接与分组，不据平台数字猜测题名、年份或认证次序。链接所对应的题名和认证编号应在下一轮选题时补齐，题目内容尚未逐一核对。

\begin{longtable}{P{1.5cm}P{13.4cm}}
\toprule 原分组 & 原始题目入口 \\\midrule\endhead
第一题 & \url{https://www.acwing.com/problem/content/3250/}\\
第一题 & \url{https://www.acwing.com/problem/content/3260/}\\
第一题 & \url{https://www.acwing.com/problem/content/3275/}\\
第一题 & \url{https://www.acwing.com/problem/content/3290/}\\
第一题 & \url{https://www.acwing.com/problem/content/3295/}\\
第一题 & \url{https://www.acwing.com/problem/content/4009/}\\
第一题 & \url{https://www.acwing.com/problem/content/4702/}\\
第一题 & \url{https://www.acwing.com/problem/content/5019/}\\
第二题 & \url{https://www.acwing.com/problem/content/3266/}\\
第二题 & \url{https://www.acwing.com/problem/content/3286/}\\
第二题 & \url{https://www.acwing.com/problem/content/3296/}\\
第二题 & \url{https://www.acwing.com/problem/content/3415/}\\
第二题 & \url{https://www.acwing.com/problem/content/4010/}\\
第二题 & \url{https://www.acwing.com/problem/content/4458/}\\
第二题 & \url{https://www.acwing.com/problem/content/4513/}\\
第二题 & \url{https://www.acwing.com/problem/content/4703/}\\
第二题 & \url{https://www.acwing.com/problem/content/5020/}\\
第二题 & \url{https://www.acwing.com/problem/content/5723/}\\
第二题 & \url{https://oj.shumeng.tech/p/CSP202605B}\\
第二题 & \url{https://oj.shumeng.tech/p/CSP202409B}\\
第二题 & \url{https://oj.shumeng.tech/p/CSP202412B}\\
第二题 & \url{https://oj.shumeng.tech/p/CSP202503B}\\
第二题 & \url{https://oj.shumeng.tech/p/CSP202506B}\\
\bottomrule
\end{longtable}
S2 没有列出第三题的具体题目链接。第11--12章需要新增第三题选题表，不能直接声称已有完整的前三题真题库。

\subsection{专题题单入口}
这些链接来自 S2 的 PDF 链接注释，其中部分显示文字只有“搜索”或“STL”。索引保留链接，不预设题单内容固定不变。
\begin{longtable}{P{4cm}P{10.9cm}}
\toprule 专题 & 原始题单链接 \\\midrule\endhead
二分 前缀和 差分 双指针 & \url{https://www.luogu.com.cn/training/1048062#problems}\\
搜索与剪枝 & \url{https://www.luogu.com.cn/training/1048981}\\
STL & \url{https://www.luogu.com.cn/training/1055276#problems}\\
简单动态规划 & \url{https://www.luogu.com.cn/training/1052248}\\
贪心 & \url{https://www.luogu.com.cn/training/1052246#problems}\\
\bottomrule
\end{longtable}

\subsection{题目索引的统一字段}
每题记录：讲义编号、题名、认证编号或平台题号、原始链接、来源文件及页码、主考点、辅助考点、层级、先修知识、建议用时、讲解位置、代码验证状态。尚未确认的题名或认证编号留为“待核对”，不推测填写。

\section{授课与编写推进安排}
\subsection{标准教学路线}
建议安排 24 次课，每次 2 课时，每课时按 45 分钟计，共 48 课时、36 小时；另配课后独立练习。课时表只作为排课起点，可根据诊断结果调整。

\begin{longtable}{P{2cm}P{7cm}P{5.8cm}}
\toprule 次数 & 内容 & 阶段交付与检查 \\\midrule\endhead
1--2 & 基础诊断、小模拟与边界 & 能独立完成简单题，并列出样例外测试。\\
3--4 & 排序、容器、输入处理 & 能解释选用容器的原因与复杂度。\\
5--6 & 字符串扫描与解析 & 完成扫描器小练习，覆盖空格和分隔符。\\
7--8 & 一维二维前缀和与差分 & 正确处理边界，区分查询和批量更新。\\
9--10 & 二分、二分答案与双指针 & 能证明单调性，手动追踪搜索区间。\\
11--12 & 贪心、基础数学与第二题训练 & 写出选择依据，识别需要性质化简的题。\\
13--14 & DFS、BFS和剪枝 & 明确状态、标记和最短步数的适用条件。\\
15--16 & 线性DP与基础背包 & 完整写出DP五要素，区分循环方向。\\
17--21 & 第三题完整案例与联调 & 完成4--6个层次不同的案例或案例阶段。\\
22--24 & 综合限时训练与复盘 & 记录错误原因，课后独立重做。\\
\bottomrule
\end{longtable}

\subsection{时间紧时的压缩路线}
以基础诊断为依据，优先保留小模拟、排序和 map/string、前缀和、二分答案、基础 DP 与两种背包，再集中练第三题。削减的是选读内容与重复接口讲解，不取消边界测试、解析训练和限时复盘。已熟练的知识改成课前自测，不要求所有学生从头听同一份接口清单。

\subsection{讲义编写的交付顺序}
\begin{enumerate}
\item \textbf{第一批：基础与工具。}完成第1--4章，同时补齐题名和认证编号。交付后应能支持第一题练习及简单字符串处理。
\item \textbf{第二批：第二题算法。}完成第5--10章，先写基础模型和代表题，再补迁移训练。同步修订概念和题意错误。
\item \textbf{第三批：第三题案例。}完成第11章的4--6个完整案例。应尽早确定选题，不等全部高级算法写完才开始。
\item \textbf{第四批：训练与合订。}补全第12章和附录，检查链接、引用、代码、目录、页码与打印效果，形成学生阅读版。
\end{enumerate}

\section{统一编写规范与完成标准}
\subsection{每节固定结构}
每节依次写明：学习目标和先修知识、问题与朴素做法、关键性质、算法推导、正确性说明、复杂度、完整代码、样例演算、边界与易错点、分层练习。避免只给模板而不解释什么时候适用，也避免只有题号没有题面摘要和来源。

\subsection{代码和符号约定}
正文代码建议统一使用 C++17 兼容写法，具体提交环境在正式授课前核对。数组题优先使用统一的 1 下标区间约定，字符串与 STL 使用其自然的 0 下标和左闭右开区间；发生转换时明确标注。求和、乘积和计数先检查上界，再决定是否使用 long long 或取模。递归代码说明最大深度，不默认无限栈空间。

每段代码应能独立编译，并配输入、输出和至少一组边界测试。以图片保存的代码转写后，要实际编译运行，不以“看上去与截图相同”作为通过标准。代码使用 listings 等环境排版，避免粘贴低清晰度截图。

\subsection{完整题解的验收标准}
\begin{itemize}
\item 题名、题号、原始链接与来源页码可追溯，转载题面的差异有说明。
\item 状态、公式、代码使用一致的含义；复杂度能够覆盖目标数据范围。
\item 代表题有完整思路和实现，部分分解法明确标注可覆盖范围。
\item 至少覆盖最小规模、重复元素、极端值和适用的特殊结构；小规模可对照暴力验证。
\item 第三题案例给出函数职责和关键状态测试，展示一次实际的错误定位过程。
\item LaTeX 成品目录完整，公式和代码不越界，跨页表格有重复表头，链接可点击。
\end{itemize}

\subsection{仓库内容安排}
\texttt{main.tex} 保存本次汇总与规划的 LaTeX 主文件；\texttt{sources/} 保存7份原始资料；\texttt{data/} 保存来源清单和从PDF提取的练习链接；\texttt{dist/} 保存可阅读的编译成品。首页说明编译方式和当前完成范围。

正式扩写时，可将正文拆为 \texttt{chapters/01-basics.tex} 等分章文件，再由新的讲义主文件统一引入。本编写方案应独立保留，避免教师用的迁移表和勘误任务混入学生版正文。新增例题、代码和测试时同步更新题目索引。

\end{document}
